\section{Conclusion}
\label{sec:conclusion}

We introduced \ours, a behavioral classification framework that distinguishes three mechanisms behind LLM false outputs: confabulations from epistemic uncertainty, incentive-driven errors from social pressure (sycophancy), and systematic errors from memorized misconceptions. Evaluating 900 questions across \gptfour and \gemini, we find a consistent distribution: approximately 61\% confabulations, 18--23\% incentive-driven, and 15--21\% systematic errors. This distribution is stable across models ($p = 0.54$), suggesting it reflects fundamental properties of RLHF-trained systems. A text-based classifier confirms that these types produce partially distinguishable signatures (71.9\% accuracy versus 33.3\% chance).

The key takeaway is that roughly 1 in 5 LLM errors involves the model overriding its own knowledge under social pressure---a qualitatively more concerning failure mode than simple ignorance, and one that existing detection methods may miss.

Future work should scale to larger question sets to narrow confidence intervals, test additional forms of social pressure beyond ``I don't think,'' and combine behavioral classification with internal-state probing on open-weight models to validate whether behavioral categories correspond to distinct mechanistic signatures. We also encourage benchmark developers to adopt mechanism-level error labels, as our framework provides a practical methodology for doing so.
